\section{Quick Sort}
O QuickSort é um dos algoritmos de ordenação mais eficientes e amplamente utilizados, baseado no conceito de divisão e conquista. A ideia central do algoritmo é escolher um elemento como pivô e, em seguida, reorganizar os elementos da lista de modo que todos os menores que o pivô fiquem à sua esquerda e todos os maiores fiquem à sua direita. Esse processo, conhecido como partição, divide a lista em duas partes que podem ser ordenadas recursivamente, aplicando o mesmo procedimento.


O algoritmo Particiona, mostrado no pseudocódigo abaixo, escolhe o último elemento do intervalo como pivô.

\lstinputlisting[language=C]{funcoes-c/h_particiona.txt}
\lstinputlisting[language=C]{funcoes-c/h_quick_sort_basico_recursivo.txt}
\lstinputlisting[language=C]{funcoes-c/h_quick_sort_basico.txt}
\lstinputlisting[language=C]{funcoes-c/h_quick_sort_recursivo.txt}
\lstinputlisting[language=C]{funcoes-c/h_quick_sort.txt}
% \begin{algorithm}
%     \caption{Particiona}\label{alg-partition:cap}
%     \DontPrintSemicolon
%     \Entrada{um intervalo [$inicio$, $fim$] e um vetor $v$, com $0 \leq inicio < fim < |v|$}
%     \Resultado{devolve $j \in [inicio, fim]$ tal que $v[inicio..j-1] \leq v[j] < v[j+1..fim]$}
%     \Inicio{
%         $piv\hat{o} \gets v[fim]$\;
%         $j \gets inicio$\;
%         \Para{$i = inicio$ \Ate $\,fim - 1\,$}{
%             \Se{$v[i] \leq piv\hat{o}$}{
%                 Troca($v[j]$, $v[i]$)\;
%                 $j \gets j + 1$\;
%             }
%         }
%         $v[fim] \gets v[j]$\;
%         $v[j] \gets piv\hat{o}$\;
%         \Retorna{j}
%     }
% \end{algorithm}

Para otimizar o uso de memória no QuickSort, é possível adotar estratégias que limitam o crescimento da pilha de execução, evitando problemas como o esgotamento de espaço em cenários desfavoráveis, por exemplo, quando o vetor está ordenado de forma decrescente. Uma abordagem eficiente é sempre processar primeiro o menor dos subvetores gerados pela partição e eliminar a segunda chamada recursiva, transformando-a em uma iteração. Dessa forma, garante-se que o tamanho do subvetor no topo da pilha será, no máximo, metade do tamanho do subvetor logo abaixo, limitando a altura da pilha a $\log_2 n$ \cite[p.~89--90]{book97338b4c}.

% \begin{algorithm}
%     \caption{QuickSort}\label{alg-quick-sort:cap}
%     \DontPrintSemicolon
%     \Entrada{um intervalo [$inicio$, $fim$] e um vetor $v$, com $0 \leq inicio \leq fim + 1 \leq |v|$}
%     \Resultado{o vetor $v[inicio..fim]$ em ordem crescente}
%     \Inicio{
%         \Enqto{$inicio < fim$}{
%             $j \gets$ Particiona($inicio$, $fim$, $v$)\;
%             \eSe{$j - inicio < fim - j$}{
%                 QuickSort($inicio$, $j-1$, $v$)\;
%                 $inicio \gets j+1$\;
%             }{
%                 QuickSort($j+1$, $fim$, $v$)\;
%                 $fim \gets j-1$\;
%             }
%         }
%     }
% \end{algorithm}

\subsection{Prova de correção}
\lipsum[19]

\subsection{Complexidade de tempo e espaço}
\lipsum[19]
